\section{Boundary Layer: Model-Theoretic Definitions and Basic Theorems}

\begin{definition}[Categoricity in a cardinal]
\label{def:categoricity}
\lean{Cardinal.Categorical}
\mathlibok
% Source: Ramsey, introduction and Corollary 7.3.
% Status: mathlib-existing.
Let \(T\) be a complete first-order theory.  In the Morley setting below,
the language is countable, \(T\) has infinite models, and the cardinals
under discussion are uncountable.  Thus the usual definition, which also
requires a model of cardinality \(\kappa\), is equivalent by
Loewenheim--Skolem to the Mathlib convention: \(T\) is
\(\kappa\)-categorical if any two models of \(T\) of cardinality
\(\kappa\) are isomorphic.
\end{definition}

\begin{definition}[Embeddings, isomorphisms, and partial embeddings]
\label{def:embeddings}
% Source: Ramsey, Definition 1.1.
% Status: local-needed; embeddings and isomorphisms are Mathlib-existing,
% but this node also includes partial embeddings.
% Mathlib alignment: partial embeddings align with FirstOrder.Language.PartialEquiv
% after importing Mathlib.ModelTheory.PartialEquiv.
An \(L\)-embedding is an injective map preserving the interpretations of
all symbols of \(L\), and an isomorphism is a bijective \(L\)-embedding.
If \(A \subseteq M\) and \(B \subseteq N\), then a map \(f:A\to B\) is
a partial embedding when
\[
  f\cup\{(c^M,c^N): c\text{ is a constant symbol of }L\}
\]
is a bijection preserving all relation and function symbols of \(L\).
\end{definition}

\begin{definition}[Elementary embeddings and substructures]
\label{def:elementary-maps}
\lean{FirstOrder.Language.ElementaryEmbedding, FirstOrder.Language.Substructure.IsElementary, FirstOrder.Language.ElementarySubstructure}
\uses{def:embeddings}
% Source: Ramsey, Definition 1.2.
% Status: local-needed; Mathlib has elementary embeddings and elementary
% substructures, but not partial elementary embeddings.
Suppose \(M\) and \(N\) are \(L\)-structures.  An \(L\)-embedding
\(j:M\to N\) is an elementary embedding if, for every formula
\(\varphi\) of \(L\) and every \(a_1,\ldots,a_n\in M\),
\[
  M\models\varphi(a_1,\ldots,a_n)
  \Longleftrightarrow
  N\models\varphi(j(a_1),\ldots,j(a_n)).
\]
This definition extends to partial embeddings in the obvious way: a
partial elementary embedding is a partial embedding \(f:A\to B\) such
that, for every formula \(\varphi\) of \(L\) and every
\(c_1,\ldots,c_n\in A\),
\[
  M\models\varphi(c_1,\ldots,c_n)
  \Longleftrightarrow
  N\models\varphi(f(c_1),\ldots,f(c_n)).
\]
If the inclusion map \(i:M\to N\) is an elementary embedding, then
\(M\) is an elementary substructure, or submodel, of \(N\), written
\(M\preceq N\).  Equivalently, \(N\) is an elementary extension of \(M\).
\end{definition}

\begin{theorem}[Tarski--Vaught test]
\label{thm:tarski-vaught-test}
\lean{FirstOrder.Language.Substructure.isElementary_of_exists}
\mathlibok
% Source: Ramsey, Theorem 1.3.
Let \(M\) be a substructure of \(N\).  Then \(M \preceq N\) if and only
if every existential formula over parameters from \(M\) that has a
witness in \(N\) already has a witness in \(M\).
\end{theorem}

\begin{theorem}[Compactness]
\label{thm:compactness}
\lean{FirstOrder.Language.Theory.isSatisfiable_iff_isFinitelySatisfiable}
\mathlibok
% Source: Ramsey, Theorem 1.4.
A theory is satisfiable if and only if every finite subset of it is
satisfiable.
\end{theorem}

\begin{theorem}[Downward Loewenheim--Skolem]
\label{thm:downward-lowenheim-skolem}
\lean{FirstOrder.Language.exists_elementarySubstructure_card_eq}
\mathlibok
% Source: Ramsey, Theorem 1.5.
Let \(M\) be an \(L\)-structure and let \(X \subseteq M\).  Then \(X\)
is contained in an elementary submodel \(N \preceq M\) whose cardinality
is controlled by \(|X|+|L|+\aleph_0\).
\end{theorem}

\begin{theorem}[Elementary chains and Loewenheim--Skolem constructions]
\label{thm:loewenheim-skolem-chains}
\uses{def:elementary-maps}
% Source: Ramsey, repeated elementary-chain constructions in Sections 4--7.
% Status: local-needed.
Let \((M_i)_{i\in I}\) be an increasing elementary chain of
\(L\)-structures.  Then the union \(\bigcup_i M_i\) is an
\(L\)-structure and each \(M_i\) is elementary in the union.  Moreover,
given a theory with infinite models and a cardinal \(\kappa\) large
enough for the language, elementary chains built by repeated
Loewenheim--Skolem reductions yield models of cardinality \(\kappa\).
\end{theorem}

\begin{proof}
\uses{thm:tarski-vaught-test,thm:downward-lowenheim-skolem}
The union structure is defined by compatibility along the chain.  The
Tarski--Vaught test proves elementarity of each stage in the union,
because any existential witness over an earlier stage appears in some
later stage and can then be reflected down the chain.  The cardinal
control is obtained by applying downward Loewenheim--Skolem at successor
stages and taking unions at limit stages.
\end{proof}

\begin{theorem}[Large models with no small infinite definable sets]
\label{thm:large-model-no-small-infinite-definable-sets}
% Source: Ramsey, construction used in the proof of Theorem 5.4.
% Status: local-needed.
% The countability hypothesis can be generalized to \(|L|<\kappa\);
% this draft keeps the countable-language version used by Ramsey.
Let \(T\) be a theory in a countable language with infinite models, and
let \(\kappa\geq\aleph_1\).  Then \(T\) has a model \(M\) of
cardinality \(\kappa\) such that every parameter-definable subset of
\(M\) is either finite or has cardinality \(\kappa\).
\end{theorem}

\begin{proof}
\uses{thm:tarski-vaught-test,thm:downward-lowenheim-skolem,thm:loewenheim-skolem-chains}
Build an elementary chain of length \(\kappa\).  At each stage, work in
the constant expansion naming the parameters already chosen.  Whenever a
formula over the current stage defines an infinite set, arrange at later
stages to add new realizations of that formula.  Downward
Loewenheim--Skolem keeps intermediate stages small, and the
Tarski--Vaught test gives elementarity at unions.

Every finite parameter tuple appears at some stage.  If the definable set
over that tuple is infinite, the construction adds new realizations
cofinally often, so the final definable set has cardinality \(\kappa\).
\end{proof}

\begin{theorem}[Skolem expansion]
\label{thm:skolem-expansion}
\notready
% Source: Ramsey, Theorem 1.6.
% Status: local-needed.
% Mathlib alignment: candidate Mathlib.ModelTheory.Skolem, but Ramsey uses an iterated expansion with enough Skolem functions for all formulas.
For every \(L\)-theory \(T\), there is an expanded language \(L^*\) and
an expanded theory \(T^*\) with built-in Skolem functions.  Each model
of \(T\) admits an expansion to a model of \(T^*\), and \(L^*\) may be
chosen with \(|L^*|=|L|+\aleph_0\).
\end{theorem}

\begin{definition}[Types over parameters]
\label{def:types-over-parameters}
\uses{def:elementary-maps}
% Source: Ramsey, Definitions 2.1, 2.2, and 2.4.
% Status: mathlib-existing for complete types and realization.
% Mathlib alignment: FirstOrder.Language.Theory.CompleteType,
% FirstOrder.Language.Theory.typeOf, FirstOrder.Language.Theory.realizedTypes,
% FirstOrder.Language.Theory.typesWith.
% These require importing Mathlib.ModelTheory.Types.
For \(A \subseteq M\), an \(n\)-type over \(A\) is a set of
\(L_A\)-formulas in \(n\) free variables consistent with
\(\operatorname{Th}_A(M)\).  A complete type decides every such formula,
and \(S_n^M(A)\) is the set of complete \(n\)-types over \(A\).
For an \(L_A\)-formula \(\varphi(v)\), the associated basic open set is
\[
  [\varphi]=\{p\in S_n^M(A):\varphi\in p\}.
\]
\end{definition}

\begin{definition}[Realization, omission, and isolation]
\label{def:isolated-and-omitted-types}
\uses{def:types-over-parameters}
% Source: Ramsey, Definitions 2.2 and 2.4.
% Status: mathlib-existing for realization; isolation is stated by singletons
% or basic opens in the topology on complete types.
% Mathlib alignment: FirstOrder.Language.Theory.typeOf,
% FirstOrder.Language.Theory.realizedTypes,
% CompleteType.isOpen_typesWith.
A type \(p\) is realized by a tuple \(a\) if every formula in \(p\) holds
of \(a\), and it is omitted by \(M\) if no tuple in \(M\) realizes it.
The type \(p \in S_n^M(A)\) is isolated if some formula
\(\varphi \in p\) determines \(p\), equivalently if
\([\varphi]=\{p\}\) in the Stone topology.
\end{definition}

\begin{definition}[Omega-stability]
\label{def:omega-stability}
\uses{def:types-over-parameters}
% Source: Ramsey, Section 3.
% Status: local-needed.
A complete theory \(T\) in a countable language is
\(\omega\)-stable if, for every model \(M\models T\), every countable
parameter set \(A\subseteq M\), and every \(n\geq 1\), the space
\(S_n^M(A)\) is countable.
\end{definition}

\begin{lemma}[Isomorphisms preserve definable set cardinalities]
\label{lem:isomorphism-preserves-definable-set-cardinality}
\uses{def:embeddings,def:elementary-maps}
% Source: formalization bridge for Ramsey, Theorem 5.4.
% Status: local-needed.
Let \(f:M\cong N\), let \(\varphi(x,y)\) be a formula, and let \(a\) be
a parameter tuple from \(M\).  Then \(f\) restricts to a bijection
\[
  \varphi(M,a)\cong \varphi(N,f(a)).
\]
In particular, these definable sets have the same cardinality.
\end{lemma}

\begin{proof}
By elementarity of the isomorphism, \(m\in M\) satisfies
\(\varphi(m,a)\) exactly when \(f(m)\) satisfies
\(\varphi(f(m),f(a))\) in \(N\).  Restricting \(f\) to the definable set
therefore gives a bijection, with inverse induced by \(f^{-1}\).
\end{proof}

\begin{lemma}[Isomorphisms preserve realized type counts]
\label{lem:isomorphism-preserves-realized-type-counts}
\uses{def:types-over-parameters,def:embeddings,def:elementary-maps}
% Source: formalization bridge for Ramsey, Theorem 5.4.
% Status: local-needed.
If \(f:M\cong N\) and \(A\subseteq M\), then \(f\) induces a bijection
between the complete types over \(A\) realized in \(M\) and the complete
types over \(f(A)\) realized in \(N\).
\end{lemma}

\begin{proof}
Send the type of a tuple \(a\) over \(A\) to the type of \(f(a)\) over
the image parameter set \(f(A)\).  Elementarity makes this independent
of the chosen realization, and the inverse is obtained from \(f^{-1}\).
\end{proof}

\begin{theorem}[Omitting types]
\label{thm:omitting-types}
\uses{def:isolated-and-omitted-types}
% Source: Ramsey, Theorem 2.5.
% Status: local-needed.
% Mathlib alignment: no matching declaration found by LeanSearch or Loogle in the current Mathlib dependency.
Let \(L\) be countable, let \(T\) be an \(L\)-theory, and let \(p\) be a
non-isolated \(n\)-type over \(\emptyset\).  Then there is a countable
model of \(T\) omitting \(p\).
\end{theorem}

\begin{definition}[Prime models over parameter sets]
\label{def:prime-over}
\uses{def:elementary-maps}
% Source: Ramsey, Definition 3.1.
% Status: local-needed.
A model \(M \models T\) is prime if it elementarily embeds into every
model of \(T\).  If \(A \subseteq M\), then \(M\) is prime over \(A\) if
every partial elementary map from \(A\) into a model of \(T\) extends to
an elementary embedding of \(M\).
\end{definition}

\begin{definition}[Homogeneous models]
\label{def:homogeneous-model}
\uses{def:elementary-maps}
% Source: Ramsey, Definition 3.6.
% Status: local-needed.
Let \(\kappa\) be infinite.  A model \(M \models T\) is
\(\kappa\)-homogeneous if every partial elementary map
\(f:A\to M\), with \(|A|<\kappa\), extends over any additional element of
\(M\).  A model is homogeneous if it is \(|M|\)-homogeneous.
\end{definition}

\begin{definition}[\((\kappa,\lambda)\)-models]
\label{def:kappa-lambda-model}
\uses{def:categoricity}
% Source: Ramsey, Definition 4.1.
% Status: local-needed.
For cardinals \(\kappa>\lambda\geq\aleph_0\), a theory \(T\) has a
\((\kappa,\lambda)\)-model if some model \(M\models T\) has cardinality
\(\kappa\) and some one-variable parameter-definable set
\(\varphi(M,a)\) has cardinality \(\lambda\).
\end{definition}

\begin{definition}[Vaughtian pairs]
\label{def:vaughtian-pair}
\uses{def:elementary-maps}
% Source: Ramsey, Definition 4.2.
% Status: local-needed.
A Vaughtian pair for \(T\) is a pair \((N,M)\) of models of \(T\) such
that \(M\) elementarily embeds into \(N\), \(M\neq N\), and for some
formula \(\varphi\) with parameters from \(M\), the definable set
\(\varphi(M)\) is infinite and \(\varphi(M)=\varphi(N)\).
\end{definition}

\begin{definition}[The pair language]
\label{def:pair-language}
\uses{def:elementary-maps}
% Source: Ramsey, Observation 4.4.
% Status: local-needed.
For an elementary pair \(M\preceq N\), let \(L_0=L\cup\{U\}\), where
\(U\) is a unary predicate interpreted as \(M\).  The pair \((N,M)\) is
then treated as a single \(L_0\)-structure.  For each \(L\)-formula
\(\varphi\), the relativization \(\varphi^U\) expresses the same formula
with all quantifiers restricted to \(U\).
\end{definition}

\begin{definition}[Order indiscernibles]
\label{def:order-indiscernibles}
\notready
% Source: Ramsey, Definition 5.1.
% Status: local-needed.
Let \((I,<)\) be an ordered set.  A sequence \((x_i)_{i\in I}\) of
distinct elements of a model \(M\) is a sequence of order indiscernibles
if any two increasing finite subsequences of the same length satisfy
the same formulas in \(M\).
\end{definition}

\begin{definition}[Skolem hull]
\label{def:skolem-hull}
\uses{thm:skolem-expansion}
% Source: Ramsey, Definition 5.2.
% Status: local-needed.
If \(M\models T^*\) and \(X\subseteq M\), then \(H(X)\) is the
\(L^*\)-substructure of \(M\) generated by \(X\).  This is the Skolem
hull of \(X\).
\end{definition}

\begin{definition}[Minimal and strongly minimal formulas]
\label{def:minimality}
\notready
% Source: Ramsey, Definition 6.1.
% Status: local-needed.
Let \(M\) be an \(L\)-structure and let \(D\subseteq M^n\) be an
infinite definable set.  The set \(D\) is minimal in \(M\) if every
definable subset of \(D\) is finite or cofinite in \(D\).  A formula
defining \(D\) is minimal.  The set or formula is strongly minimal if it
is minimal in every elementary extension of \(M\).
\end{definition}

\begin{definition}[Algebraic closure, independence, bases, and dimension]
\label{def:strongly-minimal-pregeometry}
\uses{def:minimality}
% Source: Ramsey, Definitions 6.2, 6.3, and 6.4.
% Status: local-needed.
In a structure \(M\), an element is algebraic over \(A\) if it belongs
to a finite set definable with parameters from \(A\).  If \(D\) is
strongly minimal and \(A\subseteq D\), then
\[
  \operatorname{acl}_D(A)=\{b\in D : b\text{ is algebraic over }A\}.
\]
This closure operation satisfies exchange.  A subset of \(D\) is
independent if no element lies in the closure of the others.  A basis of
\(Y\subseteq D\) is an independent subset \(B\subseteq Y\) with
\(\operatorname{acl}_D(B)=\operatorname{acl}_D(Y)\), and
\(\dim(Y)\) is the cardinality of any basis.
\end{definition}

\section{Layer 4: Isolated Types, Prime Models, and Homogeneity}

\begin{lemma}[Splitting a non-isolated basic open set]
\label{lem:isolated-splitting}
\uses{def:isolated-and-omitted-types}
% Source: Ramsey, Lemma 2.6.
% Status: local-needed.
Suppose \(\varphi(v)\) is an \(L_A\)-formula such that
\([\varphi]\subseteq S_n^M(A)\) contains no isolated types.  Then there
is a formula \(\psi(v)\) such that both
\([\varphi\wedge\psi]\) and \([\varphi\wedge\neg\psi]\) are nonempty
and contain no isolated type.
\end{lemma}

\begin{proof}
If no formula splits \([\varphi]\) into two nonempty parts, then the
complete type determined by membership in the nonempty sets
\([\varphi\wedge\theta]\) is isolated by \(\varphi\), contrary to the
hypothesis.  Choosing a genuine splitter and discarding an isolated side
if necessary gives the stated formula.
\end{proof}

\begin{lemma}[Projection of an isolated type]
\label{lem:isolated-projection}
\uses{def:isolated-and-omitted-types}
% Source: Ramsey, Lemma 2.7.
% Status: local-needed.
If \(A\subseteq M\) and \((a,b)\in M^{m+n}\) realizes an isolated type
over \(A\), then \(\operatorname{tp}^M(a/A)\) is isolated.
\end{lemma}

\begin{proof}
Let \(\varphi(v,w)\) isolate \(\operatorname{tp}^M(a,b/A)\).  Then the
formula \(\exists w\,\varphi(v,w)\) isolates the projected type of
\(a\): any formula true of \(a\) is forced by \(\varphi(v,w)\), hence by
its existential projection.
\end{proof}

\begin{lemma}[Transitivity of isolation]
\label{lem:isolated-transitivity}
\uses{def:isolated-and-omitted-types}
% Source: Ramsey, Lemma 2.8.
% Status: local-needed.
Suppose \(M\models T\), \(A\subseteq B\subseteq M\), and every finite
tuple from \(B\) realizes an isolated type over \(A\).  If
\(a\in M^n\) realizes an isolated type over \(B\), then \(a\) realizes
an isolated type over \(A\).
\end{lemma}

\begin{proof}
\uses{lem:isolated-projection}
Choose a formula \(\varphi(v,b)\) isolating the type of \(a\) over
\(B\), and choose a formula \(\theta(w)\) over \(A\) isolating the type
of the parameter tuple \(b\).  Then
\(\theta(w)\wedge\varphi(v,w)\) isolates the joint type of \((a,b)\)
over \(A\).  The preceding projection lemma isolates the type of \(a\)
over \(A\).
\end{proof}

\begin{theorem}[Density of isolated types in omega-stable theories]
\label{thm:isolated-types-dense}
\uses{def:omega-stability,def:isolated-and-omitted-types}
% Source: Ramsey, Theorem 3.3.
% Status: local-needed.
Let \(T\) be a complete theory in a countable language.  If \(T\) is
\(\omega\)-stable, then for every \(M\models T\) and every
\(A\subseteq M\), the isolated types are dense in \(S_n^M(A)\).
\end{theorem}

\begin{proof}
\uses{lem:isolated-splitting}
If some basic open set contained no isolated type, Lemma 2.6 recursively
builds a full binary tree of formulas.  Distinct branches determine
pairwise inconsistent complete types over the countable set of
parameters appearing in the tree.  This gives \(2^{\aleph_0}\) many
types over a countable parameter set, contradicting \(\omega\)-stability.
\end{proof}

\begin{lemma}[Type spaces and partial elementary maps]
\label{lem:type-space-homeomorphism}
\uses{def:types-over-parameters,def:elementary-maps}
% Source: Ramsey, Lemma 3.4.
% Status: local-needed.
If \(A\subseteq M\) and \(f:A\to N\) is partial elementary, then
\(S_n^M(A)\) is homeomorphic to \(S_n^N(f(A))\).
\end{lemma}

\begin{proof}
\uses{thm:compactness}
Send a type \(p\) over \(A\) to the type obtained by applying \(f\) to
the parameters occurring in formulas of \(p\).  Partial elementarity
preserves finite satisfiability, so compactness shows that this map is
well-defined.  Completeness of types gives injectivity, and the same
construction for \(f^{-1}\) gives the inverse.  Basic clopen sets are
carried to basic clopen sets, so the bijection is a homeomorphism.
\end{proof}

\begin{theorem}[Prime extension over a set]
\label{thm:prime-extension-over-set}
\uses{def:omega-stability,def:prime-over,def:isolated-and-omitted-types}
% Source: Ramsey, Theorem 3.5.
% Status: local-needed.
Let \(T\) be \(\omega\)-stable, let \(M\models T\), and let
\(A\subseteq M\).  There is an elementary substructure \(M_0\preceq M\)
which is prime over \(A\), and \(M_0\) may be chosen so that every
element of \(M_0\) realizes an isolated type over \(A\).
\end{theorem}

\begin{proof}
\uses{thm:tarski-vaught-test,thm:isolated-types-dense,lem:type-space-homeomorphism,lem:isolated-transitivity}
Build an increasing sequence of subsets of \(M\), starting with \(A\),
and at each successor stage add an element whose type over the current
set is isolated, if such an element exists.  At the first stage where no
new isolated type is realized outside the union, take the generated
substructure \(M_0\).

The Tarski--Vaught test and density of isolated types show that \(M_0\)
is elementary in \(M\).  To prove primeness, extend any partial
elementary map on \(A\) by transfinite induction along the construction.
At successor stages, the homeomorphism of type spaces carries the
isolating formula to an isolating formula over the image, allowing the
map to be extended elementarily.  The isolation-transitivity lemma gives
isolation over the original parameter set \(A\).
\end{proof}

\begin{theorem}[Back-and-forth for countable homogeneous models]
\label{thm:homogeneous-back-and-forth}
\uses{def:homogeneous-model,def:types-over-parameters}
% Source: Ramsey, Theorem 3.7.
% Status: local-needed.
Let \(T\) be a complete theory in a countable language.  If \(M\) and
\(N\) are countable homogeneous models of \(T\) realizing the same
complete \(n\)-types over \(\emptyset\) for every \(n\geq 1\), then
\(M\cong N\).
\end{theorem}

\begin{proof}
Enumerate \(M\) and \(N\).  A standard back-and-forth construction builds
an increasing sequence of finite partial elementary maps, alternately
putting the next element of \(M\) into the domain and the next element
of \(N\) into the range.  Homogeneity supplies the required one-point
extensions.  The union is a total elementary bijection.
\end{proof}

\section{Layer 3: Vaughtian Pairs and Few-Type Models}

\begin{lemma}[\((\kappa,\lambda)\)-models give Vaughtian pairs]
\label{lem:kappa-lambda-model-gives-vaughtian-pair}
\uses{def:kappa-lambda-model,def:vaughtian-pair}
% Source: Ramsey, Lemma 4.3.
% Status: local-needed.
If \(T\) has a \((\kappa,\lambda)\)-model with
\(\kappa>\lambda\geq\aleph_0\), then \(T\) has a Vaughtian pair.
\end{lemma}

\begin{proof}
\uses{thm:downward-lowenheim-skolem}
Let \(N\models T\) have size \(\kappa\), and let
\(X=\varphi(N,a)\) have size \(\lambda\).  By downward
Loewenheim--Skolem, choose \(M\preceq N\) of size \(\lambda\)
containing \(X\) and the parameter tuple \(a\).  Then
\(\varphi(M,a)=\varphi(N,a)\), and the pair \((N,M)\) is Vaughtian.
\end{proof}

\begin{corollary}[No small infinite definable sets without Vaughtian pairs]
\label{cor:no-small-infinite-definable-sets-without-vaughtian-pairs}
\uses{def:kappa-lambda-model,def:vaughtian-pair}
% Source: contrapositive of Ramsey, Lemma 4.3.
% Status: local-needed.
Suppose \(T\) has no Vaughtian pairs, \(M\models T\) has cardinality
\(\kappa\), and \(X\subseteq M\) is an infinite parameter-definable set.
Then \(X\) has cardinality \(\kappa\).  Equivalently, there is no
cardinal \(\lambda\) with
\[
  \aleph_0\leq\lambda<\kappa
\]
such that \(X\) has cardinality \(\lambda\).
\end{corollary}

\begin{proof}
\uses{lem:kappa-lambda-model-gives-vaughtian-pair}
If \(X\) had cardinality \(\lambda\) with
\(\aleph_0\leq\lambda<\kappa\), then \(M\) would witness a
\((\kappa,\lambda)\)-model.  Lemma 4.3 would give a Vaughtian pair,
contrary to the hypothesis.
\end{proof}

\begin{lemma}[Countable Vaughtian pair]
\label{lem:countable-vaughtian-pair}
\uses{def:vaughtian-pair,def:pair-language}
% Source: Ramsey, Lemma 4.5.
% Status: local-needed.
If \(T\) is a theory in a countable language and \(T\) has a Vaughtian
pair, then it has a Vaughtian pair \((N_0,M_0)\) with \(N_0\) countable.
\end{lemma}

\begin{proof}
\uses{thm:downward-lowenheim-skolem}
Work in the pair language \(L_0=L\cup\{U\}\), keeping the finitely many
parameters used by the witnessing formula.  Downward
Loewenheim--Skolem gives a countable elementary subpair.  Elementarity in
the pair language preserves the elementary-submodel relation and the
sentences saying that the witnessing definable set is infinite and has
no realizations outside \(U\).
\end{proof}

\begin{lemma}[Realizing a type on the submodel side]
\label{lem:realize-type-in-submodel-side}
\uses{def:pair-language,def:isolated-and-omitted-types}
% Source: Ramsey, Lemma 4.6.
% Status: local-needed.
Let \(M_0\preceq N_0\) be countable models of \(T\).  If
a type over parameters from \(M_0\) is realized in \(N_0\), then some
elementary extension of the pair \((N_0,M_0)\), viewed as an
\(L_0\)-structure, realizes that type on the \(M\)-side.
\end{lemma}

\begin{proof}
\uses{thm:compactness}
Relativize the formulas of the type to the predicate \(U\), and add the
complete diagram of the original pair.  Every finite fragment is
satisfiable because the type is realized in \(N_0\) and \(M_0\preceq
N_0\).  Compactness yields the desired elementary extension of pairs.
\end{proof}

\begin{lemma}[Realizing a type in an elementary extension of pairs]
\label{lem:realize-type-in-extension-pair}
\uses{def:pair-language,def:isolated-and-omitted-types}
% Source: Ramsey, Lemma 4.7.
% Status: local-needed.
Let \(M_0\preceq N_0\) be countable models of \(T\).  If a type over
parameters from \(N_0\) is given, then some countable elementary
extension of the pair \((N_0,M_0)\), viewed as an \(L_0\)-structure,
realizes that type.
\end{lemma}

\begin{proof}
\uses{thm:compactness}
Add the type to the elementary diagram of the pair.  Finite
satisfiability is witnessed in elementary extensions, and compactness
produces a countable elementary pair extension realizing the type.
\end{proof}

\begin{lemma}[Homogeneous countable Vaughtian pair]
\label{lem:homogeneous-vaughtian-pair}
\uses{def:vaughtian-pair,def:pair-language,def:homogeneous-model}
% Source: Ramsey, Lemma 4.8.
% Status: local-needed.
Given countable models \(M_0\preceq N_0\), there are countable models
\(M\preceq N\) forming an elementary extension of the original pair such
that \(M\) and \(N\) are homogeneous and realize the same complete types
over \(\emptyset\).
\end{lemma}

\begin{proof}
\uses{lem:realize-type-in-submodel-side,lem:realize-type-in-extension-pair,thm:homogeneous-back-and-forth}
Build an elementary chain of pairs.  At successive stages, first arrange
that types realized in the larger model are realized in the smaller
one, then arrange the extension properties needed for homogeneity on
the smaller and larger sides.  The union is countable, realizes the same
types on both sides, and is homogeneous on both sides.  The
back-and-forth theorem identifies the two sides up to isomorphism.
\end{proof}

\begin{theorem}[Vaughtian pairs produce an \((\aleph_1,\aleph_0)\)-model]
\label{thm:vaughtian-pair-gives-aleph-one-aleph-zero-model}
\uses{def:vaughtian-pair,def:kappa-lambda-model}
% Source: Ramsey, Theorem 4.9, isolated from Lemma 4.3.
% Status: local-needed.
If \(T\) is a theory in a countable language and \(T\) has a Vaughtian
pair, then \(T\) has an \((\aleph_1,\aleph_0)\)-model.
\end{theorem}

\begin{proof}
\uses{lem:countable-vaughtian-pair,lem:homogeneous-vaughtian-pair,thm:loewenheim-skolem-chains}
Obtain a countable homogeneous Vaughtian pair \((N,M)\) with a unary
formula \(\varphi(x,a)\) that has infinitely many realizations in \(M\)
and no new realizations in \(N\setminus M\).  Form an elementary chain
\((N_\alpha)_{\alpha<\omega_1}\) with each successor pair
\((N_{\alpha+1},N_\alpha)\) isomorphic to \((N,M)\).  The union has
cardinality \(\aleph_1\), while the realizations of \(\varphi(x,a)\)
remain countable, giving an \((\aleph_1,\aleph_0)\)-model.
\end{proof}

\begin{theorem}[Vaught's two-cardinal theorem]
\label{thm:vaught-two-cardinal}
\uses{def:kappa-lambda-model}
% Source: Ramsey, Theorem 4.9.
% Status: local-needed.
Let \(T\) be in a countable language.  If \(T\) has a
\((\kappa,\lambda)\)-model for some \(\kappa>\lambda\geq\aleph_0\), then
\(T\) has an \((\aleph_1,\aleph_0)\)-model.
\end{theorem}

\begin{proof}
\uses{lem:kappa-lambda-model-gives-vaughtian-pair,thm:vaughtian-pair-gives-aleph-one-aleph-zero-model}
The two-cardinal model gives a Vaughtian pair by Lemma 4.3.  The
Vaughtian-pair construction then gives an
\((\aleph_1,\aleph_0)\)-model.
\end{proof}

\begin{lemma}[Large definable set with no uncountable definable split]
\label{lem:omega-stable-large-definable-almost-minimal}
\uses{def:omega-stability}
% Source: Ramsey, Lemma 4.10.
% Status: local-needed.
Suppose \(T\) is \(\omega\)-stable, \(M\models T\), and
\(|M|\geq\aleph_1\).  Then there is an \(L_M\)-formula \(\varphi(v)\)
such that \(|\varphi(M)|\geq\aleph_1\), and for every \(L_M\)-formula
\(\psi(v)\), at least one of
\(\varphi(M)\cap\psi(M)\) and
\(\varphi(M)\setminus\psi(M)\) is countable.
\end{lemma}

\begin{proof}
If every uncountable definable set had an uncountable definable split,
then repeatedly splitting would produce a binary tree of formulas.  The
countable set of parameters appearing in the tree would support
\(2^{\aleph_0}\) many distinct one-types, contradicting
\(\omega\)-stability.
\end{proof}

\begin{lemma}[Proper extension realizing no new countable types]
\label{lem:proper-extension-realizes-no-new-countable-types}
\uses{def:omega-stability,def:isolated-and-omitted-types}
% Source: Ramsey, Lemma 4.11.
% Status: local-needed.
Suppose \(T\) is \(\omega\)-stable, \(M\models T\), and
\(|M|\geq\aleph_1\).  Then there is a proper elementary extension
\(N\succeq M\) such that every countable type over \(M\) realized in
\(N\) is already realized in \(M\).
\end{lemma}

\begin{proof}
\uses{lem:omega-stable-large-definable-almost-minimal,thm:prime-extension-over-set}
Choose the large formula supplied by Lemma 4.10 and form the complete
type consisting of the definable large subsets of it.  Realize this
type in an elementary extension, and then take a prime elementary
submodel over \(M\) together with that realization.  If a countable type
over \(M\) is realized in the prime extension, an isolating formula over
the added realization reduces the problem to a countable subset of the
large type, which is already realized in \(M\).
\end{proof}

\begin{theorem}[Two-cardinal transfer above \(\aleph_1\)]
\label{thm:two-cardinal-transfer}
\uses{def:omega-stability,def:kappa-lambda-model}
% Source: Ramsey, Theorem 4.12.
% Status: local-needed.
Suppose \(T\) is \(\omega\)-stable and has an
\((\aleph_1,\aleph_0)\)-model.  If \(\kappa>\aleph_1\), then \(T\) has
a \((\kappa,\aleph_0)\)-model.
\end{theorem}

\begin{proof}
\uses{lem:proper-extension-realizes-no-new-countable-types,thm:loewenheim-skolem-chains}
Let \(M\) be a model of size at least \(\aleph_1\) with a countable
definable set \(\varphi(M)\).  Repeatedly take proper elementary
extensions that realize no new countable types over the previous model.
The type saying ``\(\varphi(v)\) and \(v\) is not one of the old
realizations'' is countable and omitted at each stage, so the
\(\varphi\)-set does not grow.  The union of a chain of length
\(\kappa\) is a model of size \(\kappa\) with \(\varphi\)-set countable.
\end{proof}

\begin{theorem}[Models realizing few types]
\label{thm:few-types-model}
\uses{def:types-over-parameters}
% Source: Ramsey, Theorem 5.3.
% Status: local-needed.
Let \(L\) be countable and let \(T\) be an \(L\)-theory with infinite
models.  For every \(\kappa\geq\aleph_0\), there is a model
\(M\models T\) of cardinality \(\kappa\) such that, for every
\(A\subseteq M\), the model \(M\) realizes at most
\(|A|+\aleph_0\) many complete types over \(A\).
\end{theorem}

\begin{proof}
\uses{def:order-indiscernibles,def:skolem-hull,thm:skolem-expansion}
Construct such a model as the \(L\)-reduct of the Skolem hull of a
sequence of order indiscernibles in a Skolem expansion.  For a parameter
set \(A\), only \(|A|+\aleph_0\) many indiscernibles occur in Skolem
terms naming elements of \(A\).  Tuples of remaining indiscernibles are
equivalent when they have the same positions with respect to those cuts,
and there are at most \(|A|+\aleph_0\) such equivalence classes.  Skolem
terms applied to equivalent tuples realize the same \(L\)-type over
\(A\), so the reduct realizes only \(|A|+\aleph_0\) many complete
\(L\)-types over \(A\).
\end{proof}

\begin{lemma}[Non-omega-stability produces a large model realizing many types]
\label{lem:non-omega-stable-large-model-realizing-many-types}
\uses{def:omega-stability,def:types-over-parameters}
% Source: Ramsey, Theorem 5.4.
% Status: local-needed.
If \(T\) is not \(\omega\)-stable and \(\kappa\geq\aleph_1\), then
there is a model \(M\models T\) of cardinality \(\kappa\) and a
countable parameter set \(A\subseteq M\) over which \(M\) realizes
uncountably many complete types.
\end{lemma}

\begin{proof}
\uses{thm:compactness,thm:loewenheim-skolem-chains}
By non-\(\omega\)-stability, some countable parameter set \(A\) supports
uncountably many complete types.  Choose \(\aleph_1\) of these types and
add variables asking that each chosen type be realized.  Every finite
fragment is realized in an elementary extension, so compactness gives a
model realizing all chosen types.  An elementary-chain enlargement and
Loewenheim--Skolem cardinal control then produce a model of cardinality
\(\kappa\) still containing \(A\) and realizing those uncountably many
types.
\end{proof}

\begin{theorem}[Uncountable categoricity implies omega-stability]
\label{thm:uncountable-categoricity-implies-omega-stable}
\uses{def:categoricity,def:omega-stability}
% Source: Ramsey, Theorem 5.4.
% Status: local-needed.
Let \(T\) be a complete theory in a countable language with infinite
models.  If \(T\) is \(\kappa\)-categorical for some
\(\kappa\geq\aleph_1\), then \(T\) is \(\omega\)-stable.
\end{theorem}

\begin{proof}
\uses{lem:non-omega-stable-large-model-realizing-many-types,thm:few-types-model,lem:isomorphism-preserves-realized-type-counts}
Suppose \(T\) is not \(\omega\)-stable.  The preceding lemma gives a
model \(M\models T\) of size \(\kappa\) and a countable set
\(A\subseteq M\) over which \(M\) realizes uncountably many types.
Theorem 5.3 gives another model \(N\models T\) of size \(\kappa\) which
realizes at most countably many types over each countable parameter set.
By \(\kappa\)-categoricity, \(M\cong N\).  The image of \(A\) is
countable, and isomorphisms preserve the number of realized complete
types over corresponding parameter sets, contradiction.
\end{proof}

\begin{theorem}[Uncountable categoricity rules out Vaughtian pairs]
\label{thm:uncountable-categoricity-no-vaughtian-pairs}
\uses{def:categoricity,def:omega-stability,def:vaughtian-pair}
% Source: Ramsey, Theorem 5.4.
% Status: local-needed.
Let \(T\) be a complete theory in a countable language with infinite
models.  Suppose \(T\) is \(\omega\)-stable and \(T\) is
\(\kappa\)-categorical for some \(\kappa\geq\aleph_1\).  Then \(T\) has
no Vaughtian pairs.
\end{theorem}

\begin{proof}
\uses{thm:vaughtian-pair-gives-aleph-one-aleph-zero-model,thm:two-cardinal-transfer,thm:large-model-no-small-infinite-definable-sets,lem:isomorphism-preserves-definable-set-cardinality}
Assume that \(T\) has a Vaughtian pair.  The Vaughtian-pair construction
gives an \((\aleph_1,\aleph_0)\)-model.  If
\(\kappa=\aleph_1\), use this model directly; if
\(\kappa>\aleph_1\), use the two-cardinal transfer theorem to get a
\((\kappa,\aleph_0)\)-model.  Thus there is a model \(M\models T\) of
size \(\kappa\) with an infinite definable set \(\varphi(M,a)\) of
cardinality \(\aleph_0\).

The large-model construction gives a model \(N\models T\) of size
\(\kappa\) in which every infinite parameter-definable set has
cardinality \(\kappa\).  By categoricity, choose an isomorphism
\(f:M\cong N\).  It carries \(\varphi(M,a)\) bijectively to
\(\varphi(N,f(a))\).  The latter is infinite, so in \(N\) it has
cardinality \(\kappa\), contradicting the countability of
\(\varphi(M,a)\).
\end{proof}

\begin{theorem}[Uncountable categoricity implies omega-stability and no Vaughtian pairs]
\label{thm:uncountable-categoricity-implies-omega-stable-no-vaughtian-pairs}
\uses{def:categoricity,def:omega-stability,def:vaughtian-pair}
% Source: Ramsey, Theorem 5.4.
% Status: local-needed.
Let \(T\) be a complete theory in a countable language with infinite
models.  If \(T\) is \(\kappa\)-categorical for some
\(\kappa\geq\aleph_1\), then \(T\) is \(\omega\)-stable and has no
Vaughtian pairs.
\end{theorem}

\begin{proof}
\uses{thm:uncountable-categoricity-implies-omega-stable,thm:uncountable-categoricity-no-vaughtian-pairs}
The first theorem gives \(\omega\)-stability.  Applying the second
theorem with this \(\omega\)-stability rules out Vaughtian pairs.
\end{proof}

\section{Layer 2: Strong Minimality and Dimension}

\begin{lemma}[Independent tuples in a strongly minimal set have the same type]
\label{lem:independent-tuples-same-type}
\uses{def:minimality,def:strongly-minimal-pregeometry}
% Source: Ramsey, Lemma 6.5.
% Status: local-needed.
Suppose \(M,N\models T\), and \(\varphi(v)\) is a strongly minimal
formula with parameters from \(A\subseteq M_0\), where
\(M_0\models T\) and \(M_0\preceq M,N\).  If
\(a_1,\ldots,a_n\in\varphi(M)\) and
\(b_1,\ldots,b_n\in\varphi(N)\) are independent over \(A\), then
\[
  \operatorname{tp}^M(a_1,\ldots,a_n/A)
  =
  \operatorname{tp}^N(b_1,\ldots,b_n/A).
\]
\end{lemma}

\begin{proof}
Proceed by induction on \(n\).  In the one-element case, any formula
over \(A\) true of one non-algebraic point has cofinite intersection
with the strongly minimal set, hence is true of every other
non-algebraic point in any elementary extension over \(M_0\).  The
induction step applies the same finite/cofinite argument over the
previous independent tuple, using the induction hypothesis to transfer
the relevant finite cardinality statement from \(M\) to \(N\).
\end{proof}

\begin{theorem}[Dimension classification for strongly minimal theories]
\label{thm:strongly-minimal-dimension-classification}
\uses{def:minimality,def:strongly-minimal-pregeometry}
% Source: Ramsey, Theorem 6.6.
% Status: local-needed.
Suppose \(T\) is a strongly minimal theory.  If \(M,N\models T\), then
\(M\cong N\) if and only if \(\dim(M)=\dim(N)\).  More generally, if
\(\varphi(v)\) is a strongly minimal formula with parameters from
\(A\), where \(A=\emptyset\) or \(A\subseteq M_0\preceq M,N\), then
there is a bijective partial elementary map
\[
  f:\varphi(M)\to\varphi(N)
\]
whenever the corresponding strongly minimal dimensions agree.
\end{theorem}

\begin{proof}
\uses{lem:independent-tuples-same-type}
Choose bases \(B\) and \(C\) for the relevant strongly minimal sets.  If
the dimensions agree, fix a bijection \(B\to C\).  Lemma 6.5 makes this
bijection partial elementary.  By a maximality argument, extend it to a
maximal partial elementary map between subsets of
\(\varphi(M)\) and \(\varphi(N)\).  Any omitted point is algebraic over
the current domain, so an isolating formula extends the map, contradicting
maximality.  The same argument on the range makes the map bijective.
\end{proof}

\begin{lemma}[Uncountable subsets of a strongly minimal set have full dimension]
\label{lem:strongly-minimal-uncountable-subset-full-dimension}
\uses{def:minimality,def:strongly-minimal-pregeometry}
% Source: dimension bridge used in Ramsey, Theorem 7.2.
% Status: local-needed.
Let \(D\) be a strongly minimal definable set in a countable language.
If \(Y\subseteq D\) is uncountable, then
\[
  \dim(Y)=|Y|.
\]
\end{lemma}

\begin{proof}
Algebraic closure over a finite parameter tuple is countable in a
countable language, and more generally
\[
  |\operatorname{acl}_D(A)\cap Y|\leq |A|+\aleph_0.
\]
Let \(B\) be a basis of \(Y\).  Since
\(Y\subseteq\operatorname{acl}_D(B)\), the displayed estimate gives
\[
  |Y|\leq |B|+\aleph_0.
\]
Because \(Y\) is uncountable and \(B\subseteq Y\), this implies
\(|B|=|Y|\), namely \(\dim(Y)=|Y|\).
\end{proof}

\begin{theorem}[Omega-stable theories have minimal formulas]
\label{thm:omega-stable-has-minimal-formula}
\uses{def:omega-stability,def:minimality}
% Source: Ramsey, Theorem 6.7.
% Status: local-needed.
Let \(T\) be an \(\omega\)-stable theory.  If \(M\models T\), then
there is a minimal formula in \(M\).
\end{theorem}

\begin{proof}
If no minimal formula existed, begin with \(v=v\) and recursively split
each infinite definable set into two infinite definable subsets.  The
resulting binary tree of formulas gives \(2^{\aleph_0}\) complete
one-types over the countable set of parameters appearing in the tree,
contradicting \(\omega\)-stability.
\end{proof}

\begin{lemma}[Uniform finite bound without Vaughtian pairs]
\label{lem:finite-bound-without-vaughtian-pairs}
\uses{def:vaughtian-pair}
% Source: Ramsey, Lemma 6.8.
% Status: local-needed.
Suppose \(T\) has no Vaughtian pairs, \(M\models T\), and
\(\varphi(v,w)\) has parameters from \(M\).  Then there is
\(n\in\mathbb{N}\) such that whenever \(a\in M\) and
\(|\varphi(M,a)|>n\), the set \(\varphi(M,a)\) is infinite.
\end{lemma}

\begin{proof}
\uses{thm:compactness}
If no such bound existed, choose parameters \(a_n\) for which
\(\varphi(M,a_n)\) is finite of size at least \(n\).  In the pair
language, write a type asserting that \(U\) is a proper elementary
submodel, that the chosen parameters lie in \(U\), that
\(\varphi(v,a)\) has infinitely many solutions, and that all those
solutions lie in \(U\).  Finite fragments are realized by choosing
sufficiently large \(n\).  Compactness yields a Vaughtian pair,
contradicting the hypothesis.
\end{proof}

\begin{theorem}[Minimal formulas are strongly minimal without Vaughtian pairs]
\label{thm:minimal-formula-strongly-minimal}
\uses{def:vaughtian-pair,def:minimality}
% Source: Ramsey, Theorem 6.9.
% Status: local-needed.
If \(T\) has no Vaughtian pairs, then every minimal formula is strongly
minimal.
\end{theorem}

\begin{proof}
\uses{lem:finite-bound-without-vaughtian-pairs}
Let \(\varphi\) be minimal over \(M\).  If an elementary extension had a
definable infinite coinfinite subset of \(\varphi\), then Lemma 6.8
would give a uniform finite bound witnessing infinitude for both the
subset and its complement.  By elementarity, the same first-order
condition would already hold in \(M\), contradicting minimality in
\(M\).
\end{proof}

\section{Layer 1: Baldwin--Lachlan Characterization}

\begin{lemma}[Prime over an infinite definable set]
\label{lem:prime-over-infinite-definable-set}
\uses{def:omega-stability,def:vaughtian-pair,def:prime-over}
% Source: Ramsey, Lemma 7.1.
% Status: local-needed.
If \(T\) is \(\omega\)-stable with no Vaughtian pairs, \(M\models T\),
and \(X\subseteq M^n\) is infinite and definable, then \(M\) is prime
over \(X\), and no proper elementary submodel of \(M\) contains \(X\).
\end{lemma}

\begin{proof}
\uses{thm:prime-extension-over-set}
If a proper elementary submodel \(N\preceq M\) contained \(X\), then the
defining formula would have the same infinite set of realizations in
\(N\) and \(M\), making a Vaughtian pair.  Thus no such proper submodel
exists.  The prime-extension theorem gives an elementary submodel of
\(M\) prime over \(X\); since it cannot be proper, it is all of \(M\).
\end{proof}

\begin{theorem}[Omega-stability and no Vaughtian pairs imply uncountable categoricity]
\label{thm:omega-stable-no-vaughtian-pairs-implies-uncountably-categorical}
\uses{def:categoricity,def:omega-stability,def:vaughtian-pair}
% Source: Ramsey, Theorem 7.2, reverse implication.
% Status: local-needed.
Let \(T\) be a complete theory in a countable language with infinite
models, and let \(\kappa\geq\aleph_1\).  If \(T\) is
\(\omega\)-stable and has no Vaughtian pairs, then \(T\) is
\(\kappa\)-categorical.
\end{theorem}

\begin{proof}
\uses{thm:prime-extension-over-set,thm:omega-stable-has-minimal-formula,thm:minimal-formula-strongly-minimal,cor:no-small-infinite-definable-sets-without-vaughtian-pairs,lem:strongly-minimal-uncountable-subset-full-dimension,thm:strongly-minimal-dimension-classification,lem:prime-over-infinite-definable-set}
Let \(M,N\models T\) have cardinality \(\kappa\).  Choose a prime model
\(M_0\) and work over elementary embeddings, or copied elementary
submodels, of \(M_0\) inside \(M\) and \(N\).  By \(\omega\)-stability
there is a minimal unary formula over \(M_0\), and by the absence of
Vaughtian pairs it is strongly minimal.  Write its definable set as
\(D\), defined by a unary formula \(\varphi(x)\) over \(M_0\).

The sets \(D(M)\) and \(D(N)\) are infinite parameter-definable subsets.
Since there are no Vaughtian pairs, the no-small-definable-set corollary
gives
\[
  |D(M)|=\kappa=|D(N)|.
\]
The full-dimension lemma gives
\[
  \dim D(M)=\kappa=\dim D(N).
\]
The strongly minimal dimension classification theorem therefore gives a
bijective partial elementary map \(D(M)\to D(N)\) over \(M_0\).

By Lemma 7.1, \(M\) is prime over \(D(M)\), so this partial elementary
map extends to an elementary embedding \(F:M\to N\).  The image
\(F(M)\preceq N\) contains \(D(N)\).  The second part of Lemma 7.1 says
that no proper elementary submodel of \(N\) contains \(D(N)\), so
\(F(M)=N\).  Hence \(M\cong N\).
\end{proof}

\begin{theorem}[Baldwin--Lachlan characterization used by Ramsey]
\label{thm:baldwin-lachlan-characterization}
\uses{def:categoricity,def:omega-stability,def:vaughtian-pair}
% Source: Ramsey, Theorem 7.2.
% Status: local-needed.
Let \(T\) be a complete theory in a countable language with infinite
models, and let \(\kappa\geq\aleph_1\).  Then \(T\) is
\(\kappa\)-categorical if and only if \(T\) is \(\omega\)-stable and
has no Vaughtian pairs.
\end{theorem}

\begin{proof}
\uses{thm:uncountable-categoricity-implies-omega-stable-no-vaughtian-pairs,thm:omega-stable-no-vaughtian-pairs-implies-uncountably-categorical}
The forward implication is Theorem 5.4.  The reverse implication is the
fixed-cardinal Baldwin--Lachlan theorem above.
\end{proof}

\section{Layer 0: Morley's Categoricity Theorem}

\begin{corollary}[Morley's Categoricity Theorem]
\label{cor:morley-categoricity}
\uses{def:categoricity}
% Source: Ramsey, Corollary 7.3.
% Status: local-needed.
Let \(T\) be a complete theory in a countable language with infinite
models.  Then \(T\) is \(\kappa\)-categorical for some uncountable
\(\kappa\) if and only if \(T\) is \(\lambda\)-categorical for every
uncountable \(\lambda\).
\end{corollary}

\begin{proof}
\uses{thm:baldwin-lachlan-characterization}
Assume first that \(T\) is \(\kappa\)-categorical for some uncountable
\(\kappa\).  Theorem 7.2 gives that \(T\) is \(\omega\)-stable and has
no Vaughtian pairs.  Applying Theorem 7.2 again to any uncountable
\(\lambda\) shows that \(T\) is \(\lambda\)-categorical.

The converse is immediate: if \(T\) is \(\lambda\)-categorical for every
uncountable \(\lambda\), then it is categorical in particular in
\(\aleph_1\).
\end{proof}
